\documentclass[10pt]{article}

\usepackage[utf8]{inputenc}

\usepackage[T1]{fontenc}

\usepackage{amsmath}

\usepackage{amsfonts}

\usepackage{amssymb}

\usepackage[version=4]{mhchem}

\usepackage{stmaryrd}

\usepackage{graphicx}

\usepackage[export]{adjustbox}

\graphicspath{ {./images/} }



\begin{document}

$c a$, начинающегося в точке $x_{0}$ ), то для любого функционала $F[x(\tau)]$ В. ф. и. равен



\[

\int_{C_{k}^{x_{0}}(0, T)} F[x(\tau)] d W_{0, T}^{x_{0}}

\]



Часто такие интегралы определяют по условной мере $W_{0, T}^{x_{0}, y_{0}}$, порождаемой мерой Винера на пространстве траекторий $x(t)$ из $C_{k}^{x_{0}}(0, T)$ таких, что $x(T)=y_{0}$. В. Ф. и. введен Н. Винером (N. Wiener, 1923).



Применения В. ф. и. в математич. физике связаны с известным представлением Грина функции $G(x, y)$ для диффузии уравнения



\[

\frac{\partial u}{\partial t}=\Delta u+V(x) u,

\]



где $\Delta$ - оператор Лапласа, $V(x)$ - потенциал:



\[

G(x, y)=\int \exp \left\{-\int_{0}^{T} V[x(\tau)] d \tau\right\} d W_{0, T}^{x, y} .

\]



Корректность определения В. ф. и. служит математич. обоснованием использования функциональных интегралов в квантовой механике.



Jum.: [1] К а М., Вероятность и смежные вопросы в физике, пер. с англ., М., 1965; [2] Гл и м м Д., Д жаффе А., Математические методы квантовой физики. Подход с использованием функциональных интегралов, пер. с англ., М., 1984. Р. А. Минлос. ВИНТ - см. Винтовое движение.



ВИНТОВОЕ ДВИЖЕНИЕ - движение твердого тела, слагающееся из прямолинейного поступательного движения с некоторой скоростью $\boldsymbol{v}$ и вращательного движения с некоторой угловой скоростью $\omega$ вокруг оси $a a_{1}$, параллельной направлению поступательной скорости (рис. 1). Тело, совершающее стационарное В. Д., то есть В. д., при к-ром



\begin{center}

\includegraphics[max width=\textwidth]{2023_07_08_fd762c659efaaaedc47ag-1}

\end{center}



Рис. 1. Винтовое движение тела, прямолинейно движушегося по оси $a a_{1}$ и врашающегося



\begin{center}

\includegraphics[max width=\textwidth]{2023_07_08_fd762c659efaaaedc47ag-1(1)}

\end{center}



Рис. 2. Подвижный и неподвижный аксоиды. вокруг этой оси.



направление оси $a a_{1}$ остается неизменным, называется в и н то м; ось $a a_{1}-$ о с ь в инта; расстояние, проходимое любой точкой тела, лежащей на оси $a a_{1}$, за время одного оборота, - шагом $h$ винта, величина $p=v / \omega-$ п а ра м е т ром в и н т а. Если вектор $\omega$ направлен в сторону, откуда вращение тела видно происходяшим против хода часовой стрелки, то при векторах $\boldsymbol{v}$ и $\omega$, направленных в одну сторону, винт называется п р а в ы м, а в разные стороны - ле вым.



Скорость и ускорение любой точки $M$ тела, отстоящей от оси $a a_{1}$ на расстоянии $r$, численно равны



\[

v_{M}=\sqrt{v^{2}+r^{2} \omega^{2}} ; \quad w_{M}=\sqrt{w^{2}+r^{2}\left(\varepsilon^{2}+\omega^{4}\right)},

\]



где $w=d v / d t, \varepsilon=d \omega / d t$.



Когда параметр $p$ постоянен, шаг винта $h=2 \pi v / \omega=2 \pi p$ также постоянен. В этом случае всякая точка $M$ тела, не лежащая на оси $a a_{1}$, описывает винтовую линию, касатель-



\section{BИНТ}

ная к K-рой в любой точке образует с плоскостью $y z$, перпендикулярной оси $a a_{1}$, угол $\alpha=\operatorname{arctg}(h / 2 \pi r)=\operatorname{arctg} v / \omega r$.



Любое сложное движение твердого тела слагается в общем случае из серии элементарных или мгновенных В. д. Ось мгновенного В. д. называется мгн ове н ной в и нто во й осью. В отличие от оси стационарного В. д., мгновенная винтовая ось непрерывно изменяет свое положение как по отношению к системе отсчета, в к-рой рассматривается движение тела, так и по отношению к самому телу, образуя при этом две линейчатые (соприкасающиеся по прямой линии) поверхности, называемые соответственно неподвижным и подвижным аксоидами (рис. 2). Геометрич. картину движения тела можно в общем случае получить качением с продольным проскальзыванием подвижного аксоида по неподвижному, осуществляя таким путем серию тех последовательных В. д., из к-рых слагается движение тела. С. М. Тара. ВИРИАЛЬНОЕ РАЗЛОЖЕНИЕ - представление давления неидеального газа в виде ряда по степеням плотности $N / V=v^{-1}$ :



\[

P=k T v^{-1}\left(1+B_{2}(T) v^{-1}+B_{3}(T) v^{-2}+\ldots\right)

\]



где $N$ - число молекул, $V$ - объем, $T-$ температура; иногда B. p. называют также в и ри альным уравнен ием сос то я н и я (см. Состояния уравнение). Первый член соответствует давлению идеального газа, коэффициенты $B_{2}(T)$, $B_{3}(T) \ldots-$ вириальны е коэффициенты, соответствующие учету взаимодействий молекул в группах из двух, трех и т. д. молекул, поэтому В. р. называют также гру пп о в ы м р а з л о е н и е м. (Аналогичные разложения имеют место и для других термодинамич. функций.) Обычно предполагают, что газ подчиняется классич. статистике и его молекулы взаимодействуют с помощью парного потенциала сил $U(r)$. Второй вириальный коэффициент, равный



\[

B_{2}(T)=2 \pi \int_{0}^{\infty}[1-\exp (-U(r) / k T)] r^{2} d r,

\]



позволяет получить простейшее уравнение состояния для неидеального газа.



Впервые В. р. введено из эмпирич. соображений Х. Камерлинг-Оннесом (H. Kamerlingh-Onnes, 1912). В дальнейшем В. р. получали с помощью теоремы вириала. Полное В. p. можно вывести на основе канонич. или большого канонич. распределения Гиббса при помощи группового разложения, полученного Х. Урселлом (H. Ursell, 1927) и обобщенного Дж. Майером (J. Mayer, 1937):



\[

P v / k T=1-\sum_{n \geqslant 1} n \beta_{n} / v^{n}(n+1),

\]



где $\beta_{n}$ - неприводимые (не поддающиеся упрощению) групповые интегралы, связанные с вириальными коэффициентами соотношением



\[

B_{n}=-\beta_{n-1}(n-1) / n .

\]



Для них справедливы выражения:



\[

\begin{gathered}

\beta_{1}=V^{-1} \int f_{12} d r_{1} d r_{2} ; \quad \beta_{2}=(1 / 2 ! V) \int f_{12} f_{23} f_{31} d r_{1} d r_{2} d r_{3} ; \\

\beta_{3}=(1 / 3 ! V) \int\left(3 f_{12} f_{23} f_{34} f_{41}+6 f_{12} f_{23} f_{34} f_{41} f_{13}+\right. \\

\left.+f_{12} f_{23} f_{34} f_{41} f_{13} f_{24}\right) d r_{1} d r_{2} d r_{3} d r_{4} ; \\

f_{i j}=\exp \left(-U_{i j} / k T\right)-1 .

\end{gathered}

\]



Для вычисления $\beta_{n}$ в любом порядке Дж. Майером разработана диаграммная техника.



В. р. справедливы лишь для достаточно малых плотностей, вдали от точки конденсации, когда не образуются большие комплексы взаимодействующих молекул и $\beta_{n}$ можно считать не зависящими от объема $V$.



B. р. имеет место также для невырожденных квантовых газов, то есть при достаточно малой плотности.



Лит.: [1] Ла ндау Л. Д., Л и фш и ц Е. М., Статистическая физика, ч. 1, 3 изд., М., 1976, гл. 7; [2] М ай е р Дж., Ге пперт - М ай е р М., Статистическая механика, пер. с англ., 2 изд., М., 1980; [3] Х и л л Т., Статистическая механика, пер. с англ., М., 1960, гл. 5; [4] Ис и х ара А., Статистическая физика, пер. с англ., М., 1973, гл. 5.



Д. Н. Зубарев.





\end{document}